\documentclass[11pt]{article}
\usepackage[margin=1in]{geometry}
\usepackage{amsmath, amssymb}
\usepackage{listings}
\usepackage{xcolor}
\usepackage{hyperref}

\lstset{
  language=Matlab,
  basicstyle=\ttfamily\small,
  keywordstyle=\color{blue},
  commentstyle=\color{gray},
  stringstyle=\color{purple},
  showstringspaces=false,
  frame=single,
  breaklines=true,
  columns=fullflexible,
}

\title{MAT 3510 --- Homework 1 Study Notes\\
\large Sample Variance \& Polynomial Addition}
\author{}
\date{}

\begin{document}
\maketitle

These are notes on the \emph{math} behind the two problems, plus MATLAB
mechanics that trip people up. They are meant to help you write the
functions yourself, not to be handed in.

\section*{Problem 1: Sample Variance}

\subsection*{The formula}
\[
  s^2 = \frac{1}{n-1}\sum_{i=1}^n (x_i - \bar x)^2,
  \qquad
  \bar x = \frac{1}{n}\sum_{i=1}^n x_i .
\]

\subsection*{What each piece means}
\begin{itemize}
  \item $\bar x$ is the sample mean: the average of your $n$ data points.
  \item $x_i - \bar x$ is the \emph{deviation} of point $i$ from the average.
        Some deviations are positive, some negative, and they sum to zero
        by definition of the mean.
  \item Squaring, $(x_i-\bar x)^2$, makes every term non-negative (so
        positive and negative deviations don't cancel) and penalizes large
        deviations more heavily than small ones.
  \item Summing the squared deviations gives the total spread in the data.
  \item Dividing by $n-1$ (rather than $n$) rescales that total into an
        average-ish quantity.
\end{itemize}

\subsection*{Why $n-1$ and not $n$? (Bessel's correction)}
If you knew the true population mean $\mu$, then
$\frac{1}{n}\sum (x_i-\mu)^2$ would be the natural, unbiased estimate of
the population variance. But we don't know $\mu$ --- we only have
$\bar x$, which is computed \emph{from the same sample}. It is a fact
that
\[
  \bar x = \arg\min_c \sum_{i=1}^n (x_i - c)^2,
\]
i.e.\ $\bar x$ is exactly the value of $c$ that minimizes the sum of
squared deviations. So $\sum (x_i - \bar x)^2 \le \sum (x_i - \mu)^2$
always, which means using $\bar x$ in place of $\mu$ systematically
\emph{understates} the true spread. Dividing by $n-1$ instead of $n$
exactly corrects this bias on average. This is also why MATLAB's
built-in \texttt{var} divides by $n-1$ by default --- so your
\texttt{samplevar} should agree with it.

\subsection*{MATLAB mechanics: \texttt{.\string^2} vs \texttt{\string^2}}
This is the most common bug on this problem. In MATLAB, \texttt{\string^}
is \emph{matrix power} --- it only makes sense for square matrices (or
scalars). If \texttt{x} is a vector and you write \texttt{x\string^2},
MATLAB will error (or, for a scalar, just square it, which is why the bug
is easy to miss when testing on a single number).

What you actually want is \emph{elementwise} squaring: put a dot before
the operator, \texttt{x.\string^2}. The dot family
(\texttt{.\string^}, \texttt{.*}, \texttt{./}) means ``do this operation
entry-by-entry,'' as opposed to the matrix versions of those operations.

\begin{lstlisting}
x = [2 4 4 4 5 5 7 9];
x.^2      % correct: squares each entry -> [4 16 16 16 25 25 49 81]
x^2       % ERROR for a non-square, non-scalar x
\end{lstlisting}

\subsection*{Recipe for \texttt{samplevar(x)}}
\begin{enumerate}
  \item \texttt{n = length(x)}
  \item \texttt{xbar = mean(x)}
  \item \texttt{deviations = x - xbar} \quad(MATLAB broadcasts the scalar
        \texttt{xbar} against every entry of \texttt{x} automatically)
  \item \texttt{squared = deviations.\string^2}
  \item \texttt{s2 = sum(squared) / (n - 1)}
\end{enumerate}
Check your work: \texttt{samplevar(x)} should equal MATLAB's built-in
\texttt{var(x)} on the same vector, for any \texttt{x}.

\section*{Problem 2: Adding Polynomials}

\subsection*{How a polynomial becomes a vector}
A polynomial like $3x^2 + 5x - 2$ is represented by its coefficients,
listed from the \emph{highest} degree down to the constant term:
\[
  3x^2 + 5x - 2 \quad \longleftrightarrow \quad [\,3,\ 5,\ -2\,].
\]
The key subtlety: the degree that a coefficient represents depends on
\emph{where it sits relative to the end of the vector}, not on a fixed
index. In a length-3 vector, position 1 is the $x^2$ coefficient; in a
length-4 vector, position 1 is the $x^3$ coefficient. The \emph{last}
entry is always the constant term ($x^0$), regardless of length.

\subsection*{Why you can't just add the vectors directly}
Adding two polynomials means adding the coefficients of matching
degrees:
\[
  (a_2x^2+a_1x+a_0) + (b_3x^3+b_2x^2+b_1x+b_0)
  = b_3x^3 + (a_2+b_2)x^2 + (a_1+b_1)x + (a_0+b_0).
\]
If $p=[a_2,a_1,a_0]$ has length 3 and $q=[b_3,b_2,b_1,b_0]$ has length 4,
then \texttt{p + q} in MATLAB either errors (lengths differ) or, if you
tried to force it, would misalign degree-2 terms in $p$ with degree-3
terms in $q$. You must first align them \emph{by degree}.

\subsection*{The trick: pad with leading zeros}
Treat the shorter polynomial as though it has zero coefficients for the
missing high-degree terms. Concretely, $x^2+2x+3$ is really
$0\cdot x^3 + 1\cdot x^2 + 2\cdot x + 3$, so as a length-4 vector it is
$[0,1,2,3]$. Once both vectors are padded on the \textbf{left} (the
high-degree end) to the same length, position $k$ in one vector lines
up with the same degree as position $k$ in the other, and ordinary
elementwise addition is correct.

\begin{center}
\begin{tabular}{r c c c c}
  $p = x^2+2x+3$      &   & 1 & 2 & 3 \\
  padded to length 4  & 0 & 1 & 2 & 3 \\
  $q = 5x^3+x+2$       & 5 & 0 & 1 & 2 \\
  \hline
  $p+q$               & 5 & 1 & 3 & 5 \\
\end{tabular}
\end{center}
so $p+q = 5x^3+x^2+3x+5$.

\emph{Common mistake:} padding on the right instead of the left. That
would shift every existing coefficient down by one degree instead of
inserting zero coefficients for the missing high-degree terms --- the
result would be a completely different (wrong) polynomial.

\subsection*{Recipe for \texttt{polyadd(p, q)}}
\begin{enumerate}
  \item \texttt{np = length(p)}, \texttt{nq = length(q)}
  \item \texttt{n = max(np, nq)} \quad (the target length --- the
        higher-degree input doesn't need padding)
  \item Pad the shorter one on the left:
        \texttt{pp = [zeros(1, n-np), p]},
        \texttt{qq = [zeros(1, n-nq), q]}
  \item \texttt{r = pp + qq}
\end{enumerate}
Check your work by hand on an example with $p$ and $q$ of different
lengths, e.g.\ $p=[1,2,3]$, $q=[5,0,1,2]$ should give $[5,1,3,5]$.

\section*{General MATLAB tips for this homework}
\begin{itemize}
  \item \textbf{File name = function name.} For MATLAB to recognize
        \texttt{samplevar} as a callable function, it must live in a
        file named exactly \texttt{samplevar.m} (same for
        \texttt{polyadd.m}). A file like \texttt{hw1.m} containing
        \texttt{function s2 = samplevar(x)} will confuse MATLAB about
        what the callable name is.
  \item \textbf{Test against a known answer.} For Problem 1, compare
        against the built-in \texttt{var}. For Problem 2, pick small
        examples you can add by hand.
  \item \textbf{Vector vs.\ scalar operators.} Whenever you're doing
        an operation ``to every entry of a list'' (square, multiply,
        divide), reach for the dot operators \texttt{.\string^},
        \texttt{.*}, \texttt{./} rather than their plain counterparts,
        unless you specifically mean linear algebra (matrix
        multiplication, matrix powers).
\end{itemize}

\end{document}
